% =====================================================================
%  1.11 Controle de qualidade
% =====================================================================
\section{Controle de qualidade}
\label{sec:monitoramento}

Cada produto percorre a mesma cadeia de comprovação: é executado, registrado,
documentado por uma evidência, verificado e validado antes de entrar no
sistema e no reporte ao Banco. Depois da entrega, o resultado se mantém pela
operação da entidade gestora, pela tarifa e pelo suporte técnico do IAT
(\autoref{fig:evidencias}). O \autoref{qua:comprovacao} detalha a
comprovação de cada produto.

\begin{figure}[htbp]
  \centering
  \caption{Cadeia de comprovação e sustentabilidade dos resultados}
  \label{fig:evidencias}
  \fluxograma{fluxograma-evidencias}
  \fonte{Elaborado pelo IAT (2026) com base em \citeonline{mop2026}.}
\end{figure}


As informações seguem o fluxo da \autoref{fig:monitoramento}: o IAT consolida
relatórios semestrais no SIGARH e no GeoPR, o IAT e a Sanepar reportam os
indicadores a cada ano, e a UGP valida e envia ao Banco Mundial. Os
indicadores da ação estão no \autoref{qua:indicadores}, e as metas anuais, na
\autoref{tab:metas-anuais} (\autoref{sec:cronograma}).

\begin{figure}[htbp]
  \centering
  \caption{Fluxo de monitoramento, validação e reporte}
  \label{fig:monitoramento}
  \fluxograma{fluxograma-monitoramento}
  \fonte{Elaborado pelo IAT (2026) com base em \citeonline{mop2026}.}
\end{figure}

\begin{quadro}[htbp]
\caption{Indicadores do Subcomponente 3.2.b}
\label{qua:indicadores}
\footnotesize
\renewcommand{\arraystretch}{1.25}
\begin{tabularx}{\textwidth}{|L|L|Q{1.6cm}|P{2.2cm}|Q{1.4cm}|}
\hline
\cab{Indicador} & \cab{Fórmula de cálculo} & \cab{Horizonte} &
\cab{Unidade} & \cab{Meta final} \\ \hline
Poços tubulares com sistemas de abastecimento implantados &
Quantidade de projetos executados & 2028 a 2033 & Projetos (nº) & 83 \\ \hline
Domicílios atendidos com abastecimento de água &
Soma dos domicílios das comunidades contempladas & 2028 a 2033 & Domicílios & 5.810 \\ \hline
Domicílios atendidos com esgotamento -- comunidades &
70\% dos domicílios atendidos com água & 2028 a 2033 & Domicílios & 4.067 \\ \hline
Domicílios atendidos com esgotamento -- área do IDR-Paraná &
Domicílios rurais (Censo 2022) nos mananciais & 2029 a 2033 & Domicílios & 2.300 \\ \hline
Sistemas implantados pelo Programa Sanepar Rural &
Quantidade de sistemas implantados & 2027 a 2031 & Sistemas (nº) & 222 \\ \hline
População rural atendida com água e esgoto &
Soma dos habitantes beneficiados & 2028 a 2033 & Pessoas (nº) & 22.708 \\ \hline
\end{tabularx}
\fonte{Adaptado de \citeonline[Quadro 27]{mop2026} e \citeonline{custo2026}.}
\end{quadro}

\begin{landscape}
\begin{quadro}[p]
\caption{Matriz de comprovação por produto}
\label{qua:comprovacao}
\scriptsize
\renewcommand{\arraystretch}{1.3}
\begin{tabularx}{\linewidth}{|P{3.2cm}|L|L|P{3.6cm}|L|}
\hline
\cab{Produto} & \cab{Registro da execução} & \cab{Evidência de comprovação} &
\cab{Verificação e validação} & \cab{Sustentabilidade do resultado} \\ \hline
Modelo de gestão implantado & Contratos de pré-adesão; registro das associações &
  Tarifa instituída; associação regularizada & SECID (validação); IAT &
  Entidade gestora operando com tarifa e manuais de desempenho \\ \hline
Trabalho técnico-social & Cadastro domiciliar; listas de presença &
  Relatório final; termos de adesão; baixa no Termo de Cooperação & IAT e IDR-Paraná &
  Famílias comprometidas com o uso e a manutenção \\ \hline
Poço perfurado & Estudo hidrogeológico; teste de vazão; análise físico-química &
  Atesto de recebimento; laudos de vistoria; outorga & IAT (fiscalização) &
  Outorga vigente; vazão e qualidade conhecidas \\ \hline
Sistema coletivo de água & Relatórios de supervisão e de fiscalização; \textit{as built} &
  Termo de Recebimento Definitivo (TRD) & IAT; UGP e Banco Mundial no encerramento &
  Operadores capacitados; suporte técnico; avaliação semestral \\ \hline
Sistema individual de esgoto & Relatórios de medição; projeto da propriedade &
  Relatório técnico aprovado; assinatura do proprietário; teste de estanqueidade;
  comprovante de pagamento & IAT (comunidades); IAT e IDR-Paraná (mananciais);
  auditoria e amostragem \textit{in loco} & Proprietário capacitado em uso e
  manutenção preventiva \\ \hline
Sistema do Sanepar Rural & Relatórios da Sanepar & Relatório técnico aprovado &
  Sanepar & Associação de moradores; monitoramento pós-obra da Sanepar \\ \hline
Capacitação & Listas de presença & Certificados & IAT &
  Pessoas capacitadas na comunidade \\ \hline
\end{tabularx}
\fonte{Elaborado pelo IAT (2026) com base em \citeonline[Quadros 28, 29 e 31]{mop2026}.}
\end{quadro}
\end{landscape}

\pendencia{As metas progressivas do Componente~3 (Quadro 23 do MOP) precisam
ser atualizadas a partir da planilha de custos: hoje preveem esgotamento em
2027 e terminam em 2032. O MOP também calcula os domicílios com água como
``soma das comunidades'' e remete a um ``Quadro XX''.}
